\section*{Chapitre \ref{ch:iv1:timed-test-5} --- Épreuve chronométrée 5}
\markboth{\protect\omTocMark{}{Réponses~: Épreuve chronométrée 5}}{\protect\omTocMark{}{Réponses~: Épreuve chronométrée 5}}

\enlargethispage{-2\baselineskip}%
\begin{solution}{iq:iv1:timed-test-5:1}
L'espérance vaut $1.5$ heure, la variance vaut $0.75$, et la probabilité vaut $5e^{-2} \approx 0.677$.

Les intervalles entre arrivées sont des exponentielles indépendantes, d'espérance $\frac12$ et de variance $\frac14$. La troisième arrivée est la somme de trois d'entre elles, soit une variable gamma\index{loi gamma}. D'où une espérance de $\frac32$ et une variance de $\frac34$.

La troisième arrivée intervient après une heure exactement lorsqu'au plus 2 événements se produisent pendant cette heure\index{processus de Poisson}. Le nombre d'événements suit une loi de Poisson d'espérance 2, donc
\[
P = e^{-2}\left(1 + 2 + \frac{2^2}{2}\right) = 5e^{-2}.
\]
Avec $e^{-2} \approx 0.135$, on obtient $0.677$. L'espérance de $1.5$ rend ce résultat plausible. Une simulation le confirme.

\iqwrong{Utiliser une approximation normale à 3 termes, qui donne environ $0.72$. La loi gamma est dissymétrique, et le calcul exact prend une seule ligne.}
\end{solution}

\begin{solution}{iq:iv1:timed-test-5:2}
L'espérance vaut 7, et la variance vaut $30\frac13$.

$N$ suit une loi géométrique d'espérance 2 et de variance 2. Un dé a pour espérance $3.5$ et pour variance $\frac{35}{12}$.

La loi de l'espérance totale donne $2 \times 3.5 = 7$\index{loi de la variance totale}.

La loi de la variance totale donne
\[
\operatorname{Var}S = \mathbb{E}N\operatorname{Var}X + \operatorname{Var}N(\mathbb{E}X)^2 = 2 \cdot \frac{35}{12} + 2 \cdot 12.25 = 5.83 + 24.5 = 30.33.
\]
L'essentiel de la variance vient de l'ignorance du nombre de dés. Une simulation renvoie $7.0$ et $30.3$.

\iqwrong{Utiliser $\mathbb{E}N\operatorname{Var}X$ seul, ou $\operatorname{Var}(N)\operatorname{Var}(X)$. Le second terme fait intervenir le carré de l'espérance de $X$, et non sa variance.}
\end{solution}

\begin{solution}{iq:iv1:timed-test-5:3}
L'espérance a posteriori vaut $\frac23$, et la probabilité d'obtenir deux piles vaut $\frac{10}{21}$.

La loi uniforme a priori est une Beta$(1, 1)$. Après 3 piles et 1 face, la loi a posteriori est une Beta$(4, 2)$\index{loi bêta-binomiale}. Son espérance vaut $\frac46$.

Les deux lancers suivants utilisent la loi a posteriori de la loi a posteriori. La première pile a pour probabilité $\frac46$. Après elle, la loi est une Beta$(5, 2)$, si bien que la pile suivante a pour probabilité $\frac57$. Le produit vaut $\frac{20}{42} = \frac{10}{21} = 0.476$.

Remarquons que ce résultat dépasse $(\frac23)^2 = 0.444$. L'incertitude sur le biais de la pièce rend les deux lancers positivement corrélés. Une intégrale numérique de $p^2$ sur la loi a posteriori donne la même valeur.

\iqwrong{Élever au carré l'espérance a posteriori, ce qui suppose que le biais est désormais connu exactement. L'écart est la variance a posteriori de $p$.}
\end{solution}

\begin{solution}{iq:iv1:timed-test-5:4}
Le nombre espéré de suites vaut 6, et la probabilité d'obtenir exactement deux suites vaut $\frac{1}{126}$.

Le nombre de suites vaut 1 plus le nombre de paires voisines de couleurs différentes\index{variables indicatrices}. Prenons l'une quelconque des 9 paires voisines. La probabilité qu'elle soit de couleurs différentes vaut $\frac{5}{10}\cdot\frac{5}{9} \times 2 = \frac59$. D'où une espérance de $1 + 9 \cdot \frac59 = 6$.

Obtenir exactement deux suites signifie toutes les rouges puis toutes les bleues, ou l'inverse. Cela fait 2 ordres sur $\binom{10}{5} = 252$, soit $\frac{1}{126}$.

Une énumération des 252 ordres confirme ces deux résultats. C'est plus que les $5.5$ suites de 10 lancers de pièces équilibrées indépendants, car le tirage sans remise rend un changement de couleur un peu plus probable.

\iqwrong{Utiliser $\frac12$ pour chaque paire, comme pour des lancers indépendants, ce qui donne $5.5$. Une fois une rouge tirée, seules 4 des 9 boules restantes sont rouges.}
\end{solution}

\begin{solution}{iq:iv1:timed-test-5:5}
La probabilité vaut $\frac{133}{216} \approx 0.616$.

Le produit est divisible par 6 lorsqu'il possède un facteur 2 et un facteur 3. Utilisons le complémentaire et l'inclusion--exclusion\index{inclusion-exclusion}.

Aucun facteur 2 signifie que tous les dés sont impairs~: $(\frac12)^3 = \frac{27}{216}$. Aucun facteur 3 signifie qu'aucun dé ne montre 3 ni 6~: $(\frac23)^3 = \frac{64}{216}$. Aucun des deux facteurs signifie que tous les dés sont dans $\{1, 5\}$~: $(\frac13)^3 = \frac{8}{216}$.

La probabilité d'échec vaut donc $\frac{27 + 64 - 8}{216} = \frac{83}{216}$, et la réponse est $\frac{133}{216}$. Une énumération confirme le décompte de 133.

\iqwrong{Multiplier les probabilités d'avoir un facteur 2 et d'avoir un facteur 3, comme si elles étaient indépendantes. Les deux dépendent des mêmes dés~: un 6 apporte les deux.}
\end{solution}

\begin{solution}{iq:iv1:timed-test-5:6}
La probabilité vaut $1 - e^{-1/2} \approx 0.393$, et le rayon médian vaut $\sqrt{2\ln 2} \approx 1.18$.

$R^2 = X^2 + Y^2$ suit une loi exponentielle d'espérance 2. En coordonnées polaires\index{loi du khi-deux}, la densité vaut $\frac{1}{2\pi}e^{-r^2/2}r\,dr\,d\theta$, et donc $P(R > r) = e^{-r^2/2}$.

Pour $r = 1$~: $1 - e^{-1/2} = 0.393$. Pour la médiane~: $e^{-r^2/2} = \frac12$ donne $r = \sqrt{2\ln 2}$.

Comparons à une dimension~: une normale centrée réduite se trouve dans $\pm 1$ avec la probabilité $0.683$. Passer à deux dimensions fait tomber cette valeur à $0.393$. En grande dimension, la masse se déplace plus loin, un fait qui compte pour les méthodes fondées sur les distances. Une simulation le confirme.

\iqfollow{Quel est le rayon moyen~?}{$\sqrt{\pi/2} = 1.25$, l'espérance d'une variable de Rayleigh.}
\end{solution}

\begin{solution}{iq:iv1:timed-test-5:7}
La probabilité vaut $\frac{7 - 3}{7 + 3} = \frac25$.

Il s'agit du théorème du scrutin\index{théorème du scrutin}. Comptons les ordres défavorables. Si le premier vote est pour B, l'ordre est défavorable, et il existe $\binom92$ ordres de ce type. Si le premier vote est pour A et que le décompte est égal à un certain moment, réfléchissons la partie précédant la première égalité~: cela échange les votes de A et de B et donne une suite commençant par B. Ce second groupe d'ordres défavorables est donc en bijection avec le premier.

Les ordres défavorables sont au nombre de $2\binom92 = 72$ sur $\binom{10}{3} = 120$. Les ordres favorables sont donc $\frac{48}{120} = \frac25$. Une énumération le confirme.

Aux limites, un décompte de 3 à 3 donne 0, puisqu'une égalité finale s'impose, et un décompte de 7 à 0 donne 1.

\iqwrong{Utiliser la probabilité que l'avance de A soit positive à la fin, qui vaut 1 ici, ou simplement la part de A dans les votes, $\frac{7}{10}$.}
\end{solution}

\begin{solution}{iq:iv1:timed-test-5:8}
L'espérance vaut 15, et la variance vaut $52.5$.

Il existe $\binom{10}{3} = 120$ triangles possibles, chacun présent avec la probabilité $\frac18$. D'où une espérance de 15\index{variables indicatrices}.

Chaque indicatrice a pour variance $\frac18 \cdot \frac78 = \frac{7}{64}$, soit un total de $120 \times \frac{7}{64} = 13.125$.

Covariances~: deux triangles sans arête commune sont indépendants. Deux triangles partageant exactement une arête comptent cinq arêtes au total, si bien que
\[
\operatorname{Cov} = 2^{-5} - 2^{-6} = \frac{1}{64}.
\]
Chaque triangle partage une arête avec $3 \times 7 = 21$ autres, une pour chaque arête et chacun des 7 autres sommets. Il y a donc $120 \times 21$ couples ordonnés, ce qui ajoute $\frac{2520}{64} = 39.375$.

Le total vaut $13.125 + 39.375 = 52.5$. Une simulation donne $15.0$ et environ $52.5$. Remarquons que la corrélation ajoute exactement trois fois la variance indépendante.

\iqwrong{S'arrêter à $13.125$, ce qui traite les triangles comme indépendants. Partager une arête crée une dépendance positive, qui rend le décompte bien plus variable.}
\end{solution}

\begin{solution}{iq:iv1:timed-test-5:9}
Il a 302 chiffres, et le chiffre de tête est 1.

$\log_{10}2^{1000} = 301.03$\index{logarithmes}. Un nombre dont le logarithme en base 10 se situe entre 301 et 302 a 302 chiffres.

La partie fractionnaire donne les premiers chiffres~: $10^{0.03}$. Comme $\ln 10 \approx 2.3$, on a $10^{0.03} = e^{0.069} \approx 1.07$. Le nombre commence donc par 1.07, et son premier chiffre est 1.

Un calcul exact avec de grands entiers donne 302 chiffres, commençant par 1071. Remarquons qu'une petite erreur sur la cinquième décimale du logarithme, multipliée par 1000, décale la partie fractionnaire d'un centième~; c'est pourquoi les chiffres donnés dans l'énoncé suffisent pour le premier chiffre, et tout juste pour le second.

\iqfollow{À long terme, à quelle fréquence le premier chiffre de $2^n$ vaut-il 1~?}{$\log_{10}2 = 30.1~\%$. Les parties fractionnaires de $n\log_{10}2$ sont équiréparties, ce qui donne la loi de Benford pour cette suite.}
\end{solution}

\begin{solution}{iq:iv1:timed-test-5:10}
L'attente moyenne vaut $\frac{35}{9} \approx 3.89$ minutes.

Les attentes sont $U \sim \text{uniforme}(0, 10)$ et $V \sim \text{uniforme}(0, 15)$. Alors\index{somme des queues pour l'espérance}
\[
\mathbb{E}\min(U, V) = \int_0^{10}\left(1 - \frac{t}{10}\right)\left(1 - \frac{t}{15}\right)dt.
\]
En développant l'intégrande en $1 - \frac{t}{6} + \frac{t^2}{150}$, l'intégrale vaut $10 - \frac{100}{12} + \frac{1000}{450} = 10 - 8.33 + 2.22 = 3.89$. Une simulation renvoie $3.89$.

Vérifications~: l'attente doit être inférieure à 5, l'attente moyenne de A seul. La règle pour des lignes de Poisson de mêmes fréquences donne $1/(\frac{1}{10} + \frac{1}{15}) = 6$, nettement plus long~: des bus réguliers sont préférables à des bus aléatoires.

\iqwrong{Répondre 3, en combinant les deux attentes moyennes comme des taux en parallèle, $\frac15 + \frac{1}{7.5}$. Les attentes sont uniformes, et non exponentielles.}
\end{solution}

\begin{solution}{iq:iv1:timed-test-5:11}
Changer de porte gagne avec probabilité $\frac{4}{15}$, contre $\frac15$ pour rester.

Votre porte cache le lot avec probabilité $\frac15$, et l'action de l'animateur ne change pas cela~: l'animateur peut toujours ouvrir une porte à chèvre. Les quatre autres portes détiennent donc le lot avec probabilité $\frac45$\index{formule de Bayes}, et l'animateur a montré qu'une d'entre elles ne le cache pas. Ces $\frac45$ se répartissent désormais sur 3 portes.

Chaque porte fermée que vous n'avez pas choisie a donc une probabilité $\frac{4}{15}$. La somme est correcte~: $\frac15 + 3 \times \frac{4}{15} = 1$.

Une simulation avec un animateur au hasard donne $0.267$. Changer de porte est meilleur d'un tiers. Avec $n$ portes et une seule ouverte, la probabilité vaut $(n - 1)/(n(n - 2))$. Pour $n = 3$, cela donne $\frac23$, le cas bien connu.

\iqwrong{Répondre $\frac14$, en traitant les quatre portes fermées comme équiprobables. Votre porte n'est pas comme les autres~: l'animateur n'avait pas le droit de l'ouvrir, donc son choix ne dit rien à son sujet.}
\end{solution}

\begin{solution}{iq:iv1:timed-test-5:12}
Hoeffding donne $e^{-8} \approx 3.4 \times 10^{-4}$, Chebyshev donne $0.0625$ pour l'événement bilatéral, et l'approximation normale donne environ $5 \times 10^{-5}$.

Ici $t = 20$ et $n = 100$, donc $2t^2/n = 8$\index{inégalité de Hoeffding}. La borne vaut $e^{-8}$.

Chebyshev utilise la variance de 25~: $P(|X - 50| \ge 20) \le 25/400 = 0.0625$. Une seule queue vaut au plus la moitié de cette valeur si l'on utilise la symétrie, soit $0.031$.

L'approximation normale avec correction de continuité utilise $z = 19.5/5 = 3.9$, avec une queue de $4.8 \times 10^{-5}$. La queue binomiale exacte vaut $3.9 \times 10^{-5}$.

Chebyshev est donc trop grossière d'un facteur mille, et Hoeffding d'un facteur d'environ 9, ce qui est bon pour une borne qui ne suppose rien d'autre que le caractère borné. Chacune est valide~: elles échangent seulement la finesse contre la généralité.

\iqwrong{Prendre une borne pour une estimation. Une borne supérieure dit ce qui ne peut pas arriver, et elle peut être très éloignée de la valeur réelle.}
\end{solution}
